% ch2_d 第2章 §2.2–§2.3 (书页 47–57 / p062–p072)
然而，对于闭合路径，我们也许无法为路径上所有的态挑选出一个共同的相位。这是因为，对于闭合路径，$|\pmb{n}(T)\rangle$ 与 $|\pmb{n}(0)\rangle$ 表示同一个自旋态，而由平行输运所确定的二者的相位选择可能并不一致。这一偏差恰好就是闭合路径的贝里相。选取自旋态的不同相位，即 $|\pmb{n}(t)\rangle \to \mathrm{e}^{\mathrm{i}\phi(t)}|\pmb{n}(t)\rangle$，会引起贝里相的改变 $\theta_B \to \theta_B + \phi(0) - \phi(T)$；由此我们看到，对于 $\pmb{n}(0) = \pmb{n}(T)$ 的闭合路径，贝里相是良定义的（在相差 $2\pi$ 的整数倍的意义下），因为 $\mathrm{e}^{\mathrm{i}\phi(0)} = \mathrm{e}^{\mathrm{i}\phi(T)}$。因此，讨论闭合路径的贝里相，或者比较具有相同起点与终点的两条路径的贝里相之差，都是有意义的。

我们看到，如果能为所有的态定义一个共同的相位，贝里相就为零。环路具有非零的贝里相，意味着不可能为环上所有的态选出一个共同的相位。贝里相的数值表示定义这种共同相位时的阻挫（frustration）程度。平行输运以及平行输运中的阻挫这些概念非常重要。电磁场与引力场都是广义的贝里相，它们描述的是某些更一般的矢量（而不是如上面所讨论的二维矢量，即复数）在平行输运中的阻挫。

贝里相是一种几何相位：只要两条路径 $\pmb{n}_1(t)$ 与 $\pmb{n}_2(t)$ 在单位球上描出相同的轨迹，它们就具有相同的贝里相。如果单位球上的轨迹 $\pmb{n}(t)$ 张出的立体角为 $\Omega$，那么该路径的贝里相就由下式给出（见问题 2.3.4）：
\begin{equation*}
\theta_B = \frac{\Omega}{2}\tag{2.3.5}
\end{equation*}

\subsection{贝里相与运动方程}

\begin{keypoints}
\item 贝里相可以影响系统的动力学性质。它可以改变运动方程。
\end{keypoints}

路径积分表示 (2.3.3) 可以推广到自旋-$S$ 系统，也可以推广到非零哈密顿量 $H = \pmb{B}\cdot\pmb{S}$ 的情形。相干态（即 $\pmb{n}\cdot\pmb{S}$ 的本征值最大（为 $S$）的本征态），记作 $|\pmb{n}\rangle$，可以写成上面讨论过的 $2S$ 个自旋-$\frac{1}{2}$ 相干态的直积：
\begin{equation*}
|\pmb{n}\rangle = |z\rangle \otimes |z\rangle \otimes \dots \otimes |z\rangle
\end{equation*}
于是，贝里相是自旋-$\frac{1}{2}$ 贝里相的 $2S$ 倍：
\begin{equation*}
\theta_B = \oint dt\, 2S\,\mathrm{i}\, z^{\dagger}\dot{z} = S\Omega\tag{2.3.6}
\end{equation*}
路径积分表示于是成为
\begin{equation*}
\mathrm{i}G(\pmb{n}_2, \pmb{n}_1, t) = \langle \pmb{n}_2 | U(t,0) | \pmb{n}_1 \rangle = \int \mathcal{D}\Big[\frac{\pmb{n}(t)}{2\pi}\Big]\, \mathrm{e}^{\mathrm{i}\int \mathrm{d}t\, [2S\,\mathrm{i}z^{\dagger}\dot{z} - \pmb{B}\cdot\pmb{n}S]}\tag{2.3.7}
\end{equation*}
让我们把
\begin{equation*}
S[\pmb{n}(t)] = \int \mathrm{d}t\, [2S\,\mathrm{i}z^{\dagger}\dot{z} - \pmb{B}\cdot\pmb{n}S]\tag{2.3.8}
\end{equation*}
当作一个经典作用量，并考虑由这样一个作用量所支配的自旋的经典运动。为得到经典运动方程，让我们比较 $S[\pmb{n}(t)]$ 与 $S[\pmb{n}(t)+\delta\pmb{n}(t)]$。（注意，路径 $\pmb{n}(t)$ 与 $\pmb{n}(t)+\delta\pmb{n}(t)$ 具有相同的起点和终点。）由贝里相 (2.3.6) 的几何解释可以发现，两条路径的贝里相之差等于 $S$ 乘以两条轨迹之间的面积：$\delta\theta_B = \int dt\, [S\pmb{n}\cdot(\dot{\pmb{n}}\times\delta\pmb{n})]$。于是，
\begin{equation*}
S[\pmb{n}(t)+\delta\pmb{n}(t)] - S[\pmb{n}(t)] = \int \mathrm{d}t\, [S\pmb{n}\cdot(\dot{\pmb{n}}\times\delta\pmb{n}) - \pmb{B}\cdot\delta\pmb{n}\,S]
\end{equation*}
这给出运动方程
\begin{equation*}
\pmb{n}\times\dot{\pmb{n}} = \pmb{B} - \pmb{n}(\pmb{n}\cdot\pmb{B})
\end{equation*}
（注意 $\delta\pmb{n}$ 总是垂直于 $\pmb{n}$），或者
\begin{equation*}
\dot{\pmb{S}} = \pmb{B}\times\pmb{S}\tag{2.3.9}
\end{equation*}
其中 $\pmb{S} = S\pmb{n}$ 对应于经典自旋矢量。这是一个奇怪的运动方程：速度（而不是加速度）正比于 $\pmb{B}$ 所代表的力。更奇怪的是，速度指向与力垂直的方向。然而，这恰好就是自旋的正确运动方程。我们看到，要恢复正确的自旋运动方程，贝里相是必不可少的。

\subsection*{问题 2.3.1}
证明图 2.4 中定义的 $\theta_B$ 实际上就是自旋-1 自旋的贝里相。

\subsection*{问题 2.3.2}
实际上，用不连续的路径 $s_z(t)$ 来描述自旋-$\frac{1}{2}$ 系统的路径积分表述也是可能的。设自旋-$\frac{1}{2}$ 系统的 $H = a\pmb{\sigma}^1 + b\pmb{\sigma}^3$，求时间演化算符 $\langle s_z'|U(t,0)|s_z\rangle$ 的路径积分表示。（提示：$\sum_{s_z=\pm 1/2}|s_z\rangle\langle s_z| = 1$。）

\subsection*{问题 2.3.3}
用旋量表示 (2.3.1) 计算闭合路径（$\phi = 0 \to 2\pi$，$\theta$ 固定）的贝里相。

对同一路径改用另一种旋量表示，计算贝里相：
\begin{equation*}
z = \begin{pmatrix} \cos\frac{\theta}{2} \\[2pt] \mathrm{e}^{-\mathrm{i}\phi}\sin\frac{\theta}{2} \end{pmatrix}
\end{equation*}
从这个例子我们看到，贝里相具有 $2\pi$ 的模糊性，依赖于所选的旋量表示。因此，贝里相 $\theta_B$ 不能直接用路径 $\pmb{n}$ 表示出来。这正是我们必须引入旋量表示来表述贝里相的原因。然而，$\mathrm{e}^{\mathrm{i}\theta_B}$ 是良定义的，并且只依赖于路径 $\pmb{n}(t)$。要得到一个良定义的路径积分，这已经足够。

\subsection*{问题 2.3.4}
对小闭合路径 $(\theta,\phi) \to (\theta+d\theta, \phi) \to (\theta+d\theta, \phi+d\phi) \to (\theta, \phi+d\phi) \to (\theta, \phi)$ 证明式 (2.3.5)。

对一般的大闭合路径证明式 (2.3.5)。

\subsection*{问题 2.3.5}
自旋与球面上的粒子：
\begin{enumerate}
\item 证明：设 $\pmb{B} = 0$，自旋作用量 (2.3.8) 可以写成
\begin{equation*}
S[\pmb{n}(t)] = \int \mathrm{d}t\, (A_{\theta}\dot{\theta} + A_{\phi}\dot{\phi})
\end{equation*}
求 $A_{\theta}$ 与 $A_{\phi}$。
\item 上述作用量可以看作下面这个作用量在 $m \to 0$ 极限下的结果：
\begin{equation*}
S_p[\pmb{n}(t)] = \int \mathrm{d}t\, \Big(A_{\theta}\dot{\theta} + A_{\phi}\dot{\phi} + \frac{1}{2}m\dot{\pmb{n}}^2\Big)
\end{equation*}
$S_p$ 描述一个质量为 $m$ 的粒子在单位球面上的运动。该粒子还经受一个由 $(A_{\phi}, A_{\theta})$ 所描述的磁场。证明 $(A_{\phi}, A_{\theta})$ 给出一个均匀磁场，并证明球面上的总磁通量子数为 $2S$。
\item 我们知道，均匀磁场中的粒子具有朗道能级结构。同一朗道能级中的所有态都具有相同的能量。各朗道能级被恒定的能隙 $\hbar\omega_c$ 分开，并且当 $m \to 0$ 时 $\omega_c \to \infty$。因此，在 $m \to 0$ 的极限下，只有第一朗道能级中的态出现在低能希尔伯特空间中。若均匀磁场的总磁通量子数为 $N_{\phi}$，求球面上粒子第一朗道能级中的态数。
\end{enumerate}

\subsection*{问题 2.3.6}
自旋波：考虑自旋-$S$ 量子自旋链 $H = \sum_i J\pmb{S}_i\cdot\pmb{S}_{i+1}$。当 $J < 0$ 时，经典基态是 $\pmb{S}_i = S\hat{z}$ 的铁磁态；当 $J > 0$ 时，是 $\pmb{S}_i = S\hat{z}(-)^i$ 的反铁磁态。
\begin{enumerate}
\item 写出自旋链的作用量。
\item 求围绕铁磁基态的小涨落 $\delta\pmb{S}_i = \pmb{S}_i - S\hat{z}$ 的运动方程。把运动方程变换到频率空间和动量空间，求出色散关系。证明对小的 $k$ 有 $\omega \propto k^2$。
\end{enumerate}

%FIG{2.5}: {双势阱。}
%FIG{2.6}: {(a) 单个瞬子解可以视为倒转势的两个极大点之间的一种运动。(b) 同一运动在时空中的作图。}
\begin{enumerate}
\setcounter{enumi}{2}
\item 求围绕反铁磁基态的小涨落 $\delta\pmb{S}_i = \pmb{S}_i - S\hat{z}(-)^i$ 的运动方程。求出色散关系。证明当 $k$ 接近 $\pi$ 时 $\omega \propto |k - \pi|$。
\end{enumerate}

有趣的是，铁磁态之上与反铁磁态之上的自旋波具有非常不同的色散。在经典统计物理中，铁磁态与反铁磁态具有相同的热力学性质；而正是贝里相使二者在量子层次上大不相同。

\section{路径积分表述的应用}

\subsection{穿越势垒的隧穿}

\begin{keypoints}
\item 有时，初始组态与最终组态由几条驻定路径相连。作用量最小的路径代表默认路径，其余路径代表瞬子效应。
\item 穿越势垒的隧穿是一种瞬子效应。
\end{keypoints}

%FIG{2.7}: {多瞬子组态——瞬子气体。}

路径积分不仅使我们能够计算诸如线性响应之类的微扰效应，还能使我们计算非微扰效应。这里我们想研究一个简单的例子，以理解路径积分如何给出非微扰的结果。我们将研究双势阱（图 2.5）中的隧穿。在没有隧穿时，粒子在某一极小点附近涨落，可用频率为 $\omega_0$ 的谐振子来描述。这些涨落给出如下传播子：
\begin{equation*}
\langle x_0|\mathrm{e}^{-TH}|x_0\rangle = \langle -x_0|\mathrm{e}^{-TH}|-x_0\rangle = \mathrm{e}^{-T\omega_0/2}, \quad \langle x_0|\mathrm{e}^{-TH}|-x_0\rangle = 0.
\end{equation*}
然而，由于隧穿，上述结果并不完全正确。有隧穿时，粒子可以在两个极小点之间来回穿行，传播子会获得一些额外的贡献。我们指出，隧穿是一种非微扰效应：围绕某一个极小点做微扰永远不可能给出隧穿。

为了研究隧穿效应，我们采用虚时路径积分和鞍点近似。除了分别连接 $x_0$ 到 $x_0$ 与 $-x_0$ 到 $-x_0$ 的两条平庸驻定路径 $x(t) = x_0$ 和 $x(t) = -x_0$ 之外，还有另一条连接 $x_0$ 与 $-x_0$ 的驻定路径。这样的路径使作用量 $\int \mathrm{d}\tau\,[\frac{1}{2}m\dot{x}^2 + V(x)]$ 取极小，并满足 $m\frac{\mathrm{d}^2 x}{\mathrm{d}\tau^2} = V'(x)$（译注：原文这两个公式分别印作含 $V'(x)$ 的作用量与 $m\frac{\mathrm{d}^2 x}{\mathrm{d}\tau^2} = V(x)$，撇号位置互易）。这是倒转势 $-V(x)$ 的牛顿方程（图 2.6(a)），其解如图 2.6(b) 所示。我们看到，在 $\pm x_0$ 之间的切换发生在一段很短的时间间隔内，这一事件称为一个瞬子（instanton）。图 2.6(b) 中的路径并不是连接 $x_0$ 与 $-x_0$ 的唯一驻定路径，图 2.7 中的多瞬子路径也是一条驻定路径。类似地，多瞬子路径还给出连接 $x_0$ 到 $x_0$ 以及 $-x_0$ 到 $-x_0$ 的其他驻定路径。

于是，为了计算传播子，我们需要对来自不同驻定路径的所有贡献求和。对于从 $x_0$ 到 $x_0$ 的传播子，我们有
\begin{align*}
\langle x_0|\mathrm{e}^{-TH}|x_0\rangle &= \mathrm{e}^{-T\omega_0/2} \sum_{n=\mathrm{even}} \int_0^T \mathrm{d}\tau_n \dots \int_0^{\tau_2} \mathrm{d}\tau_1 \left(K\mathrm{e}^{-S_0}\right)^n \\
&= \mathrm{e}^{-T\omega_0/2} \sum_{n=\mathrm{even}} \frac{T^n}{n!} \left(K\mathrm{e}^{-S_0}\right)^n = \cosh(TK\mathrm{e}^{-S_0})
\end{align*}
其中 $S_0$ 是瞬子的最小作用量，$K$ 来自围绕单个瞬子的涨落的路径积分。$K$ 应当具有频率的量纲，其数值大致由粒子撞击势垒的频率给出。类似地，我们还有
\begin{align*}
\langle x_0|\mathrm{e}^{-TH}|-x_0\rangle &= \mathrm{e}^{-T\omega_0/2} \sum_{n=\mathrm{odd}} \int_0^T \mathrm{d}\tau_n \dots \int_0^{\tau_2} \mathrm{d}\tau_1 \left(K\mathrm{e}^{-S_0}\right)^n \\
&= \mathrm{e}^{-T\omega_0/2} \sum_{n=\mathrm{odd}} \frac{T^n}{n!} \left(K\mathrm{e}^{-S_0}\right)^n = \sinh(TK\mathrm{e}^{-S_0})
\end{align*}
因此，
\begin{equation*}
\langle \psi_{\pm}|\mathrm{e}^{-TH}|\psi_{\pm}\rangle = \mathrm{e}^{-T\omega_0/2}\left[\cosh(TK\mathrm{e}^{-S_0}) \pm \sinh(TK\mathrm{e}^{-S_0})\right] = \mathrm{e}^{-T(\frac{\omega_0}{2} \pm K\mathrm{e}^{-S_0})}
\end{equation*}
其中 $|\psi_{\pm}\rangle = \frac{|x_0\rangle \pm |-x_0\rangle}{\sqrt{2}}$。我们看到 $|\psi_{\pm}\rangle$ 是能量本征态，其能量为 $E = \frac{\omega_0}{2} \pm K\mathrm{e}^{-S_0}$。

为了求出 $S_0$，注意到由能量守恒有 $\dot{x} = \sqrt{2V(x)/m}$。于是，
\begin{equation*}
S_0 = \int_{-\infty}^{+\infty} \mathrm{d}\tau\, \Big[\frac{1}{2}m\dot{x}^2 + V(x)\Big] = \int_{-x_0}^{+x_0} \mathrm{d}x\, \sqrt{2mV(x)}
\end{equation*}
恢复 $\hbar$ 后，我们得到 $\Delta E = 2K\mathrm{e}^{-\hbar^{-1}\int_{-x_0}^{+x_0} \mathrm{d}x\, \sqrt{2mV(x)}}$，这正是 WKB 结果。然而，路径积分确实教给我们一件新东西，即隧穿时间 $t_{tun}$——与单个瞬子相联系的时间尺度。对于粒子穿过一个势垒需要多长时间，我们有了一幅非常简单的图像。

\subsection*{问题 2.4.1}
设双势阱势由 $V(x) = \frac{g}{4}(x^2 - x_0^2)^2$ 给出。
\begin{enumerate}
\item 估算 $S_0$ 和隧穿时间 $t_{tun}$（即瞬子的大小）。
\item 设 $K \sim \sqrt{\frac{g x_0^2}{m}}$（这是极小点之一附近的振动频率），估算时间方向上瞬子的（平均）密度。
\end{enumerate}

%FIG{2.8}: {具有亚稳态的势。}
%FIG{2.9}: {(a) 单次反弹解可以视为倒转势中的一种运动。(b) 同一运动在时空中的作图。}
\begin{enumerate}
\setcounter{enumi}{2}
\item 若 $S_0 \gg 1$ 并且各瞬子互不重叠，半经典图像就是对隧穿的一种良好描述。问 $g$ 和 $x_0$ 取什么值时，上述两个条件都能满足？
\end{enumerate}

\subsection{亚稳态的命运}

\begin{keypoints}
\item 亚稳态的衰变也是一种瞬子效应。
\end{keypoints}

在本节中，我们考虑图 2.8 所示势 $V(x)$ 中的一个粒子。粒子初始时位于 $x_0$。我们想计算粒子穿过势垒逃逸的衰变率。为此我们计算 $\langle x_0|\mathrm{e}^{-TH}|x_0\rangle$。如果只考虑围绕 $x_0$ 的小涨落，则可以作近似 $V(x) = \frac{1}{2}m\omega_0^2(x - x_0)^2$，并对大的 $T$ 得到 $\langle x_0|\mathrm{e}^{-TH}|x_0\rangle = \mathrm{e}^{-T\omega_0/2}$。解析延拓到实时，我们看到 $\langle x_0|\mathrm{e}^{-\mathrm{i}TH}|x_0\rangle = \mathrm{e}^{-\mathrm{i}T\omega_0/2}$，没有衰变。为了找到衰变，我们需要寻找一个能使 $|\langle x_0|\mathrm{e}^{-\mathrm{i}TH}|x_0\rangle|$ 随时间减小的过程。我们将看到，隧穿可以引起衰变。在路径积分中，隧穿同样由瞬子来描述。

虚时拉格朗日量 $L = \frac{1}{2}m\dot{x}^2 + V(x)$ 有一条连接 $x_0$ 到 $x_0$ 的非平庸驻定路径 $x_c(t)$。这条驻定路径（称为一次反弹（bounce）或一个瞬子）的存在，可以从势 $-V(x)$ 中的经典运动看出来（图 2.9）。把反弹包括进来，并重复 2.4.1 节中的步骤，我们发现
\begin{align*}
\langle x_0|\mathrm{e}^{-TH}|x_0\rangle &= \mathrm{e}^{-T\omega_0/2} \sum_n \int_0^T \mathrm{d}\tau_n \int_0^{\tau_n} \mathrm{d}\tau_{n-1} \dots \int_0^{\tau_2} \mathrm{d}\tau_1 \left(K\mathrm{e}^{-S_0}\right)^n \\
&= \mathrm{e}^{-T(\frac{\omega_0}{2} - K\mathrm{e}^{-S_0})}\tag{2.4.1}
\end{align*}
其中 $S_0$ 是反弹的作用量，$K$ 是围绕反弹的涨落所产生的系数。乍看之下，反弹似乎只是把基态能量 $\frac{\omega_0}{2}$ 修正了一小量 $\Delta E = -K\mathrm{e}^{-S_0}$。由于反弹在经过很长时间后又回到 $x_0$，很难看出任何衰变。事实上我们将看到，系数 $K$ 是虚数，并且在实时 $\langle x_0|\mathrm{e}^{-\mathrm{i}TH}|x_0\rangle = \mathrm{e}^{-\mathrm{i}T\frac{\omega_0}{2} - T|K|\mathrm{e}^{-S_0}}$。反弹精确地描述了 $|x_0\rangle$ 态的衰变。这是一个相当惊人的结果。

为了理解 $K$ 为什么是虚数，我们需要讨论如何计算 $K$。这里的讨论同样适用于 2.4.1 节中隧穿振幅里的 $K$。

从我们在式 (2.4.1) 中纳入反弹的方式可以看到，$K\mathrm{e}^{-S_0}$ 是围绕 $x_c(t)$ 与围绕 $x = x_0$ 的两个路径积分之比：
\begin{equation*}
K\mathrm{e}^{-S_0} = \frac{\int_{x_c} \mathcal{D}\delta x\, \mathrm{e}^{-S}}{\int_{x=x_0} \mathcal{D}\delta x\, \mathrm{e}^{-S}}\tag{2.4.2}
\end{equation*}
对于小涨落，拉格朗日量可以展开到二阶。在 $x_c$ 附近，我们有
\begin{equation*}
S = S_0 + \int \mathrm{d}\tau\, \delta x \frac{1}{2}\Big[-m\frac{\mathrm{d}^2}{\mathrm{d}\tau^2} + V''(x_c(\tau))\Big]\delta x
\end{equation*}
而在 $x = x_0$ 附近有
\begin{equation*}
S = \int \mathrm{d}\tau\, \delta x \frac{1}{2}\Big[-m\frac{\mathrm{d}^2}{\mathrm{d}\tau^2} + V''(x_0)\Big]\delta x
\end{equation*}
于是，完成高斯积分后，我们得到
\begin{align*}
K &= \frac{\int_{x_c} \mathcal{D}\delta x\, \mathrm{e}^{-\int \mathrm{d}\tau\, \frac{1}{2}\delta x[-m\frac{\mathrm{d}^2}{\mathrm{d}\tau^2} + V''(x_c(\tau))]\delta x}}{\int_{x=x_0} \mathcal{D}\delta x\, \mathrm{e}^{-\int \mathrm{d}\tau\, \frac{1}{2}\delta x[-m\frac{\mathrm{d}^2}{\mathrm{d}\tau^2} + V''(x_0)]\delta x}} \\
&= \left(\frac{\operatorname{Det}\left(-m\frac{\mathrm{d}^2}{\mathrm{d}\tau^2} + V''(x_0)\right)}{\operatorname{Det}\left(-m\frac{\mathrm{d}^2}{\mathrm{d}\tau^2} + V''(x_c)\right)}\right)^{1/2}\tag{2.4.3}
\end{align*}
然而，式 (2.4.3) 是没有意义的，因为算符 $-m\frac{\mathrm{d}^2}{\mathrm{d}\tau^2} + V''(x_c)$ 具有一个零本征值。为了看清这一点，我们注意到 $x_c(\tau)$ 与 $x_c(\tau + \eta)$ 都是驻定路径，它们满足运动方程
\begin{equation*}
-m\frac{\mathrm{d}^2 x_c(\tau)}{\mathrm{d}\tau^2} + V'(x_c(\tau)) = -m\frac{\mathrm{d}^2 x_c(\tau+\eta)}{\mathrm{d}\tau^2} + V'(x_c(\tau+\eta)) = 0
\end{equation*}
将上式展开到 $\eta$ 的线性阶，我们发现 $\eta\dot{x}_c(\tau) = x_c(\tau+\eta) - x_c(\tau)$ 满足
\begin{equation*}
\Big[-m\frac{\mathrm{d}^2}{\mathrm{d}\tau^2} + V''(x_c(\tau))\Big]\dot{x}_c(\tau) = 0
\end{equation*}
于是，$\dot{x}_c(\tau)$ 是 $\left(-m\frac{\mathrm{d}^2}{\mathrm{d}\tau^2} + V''(x_c)\right)$ 的一个零本征值本征态。

从上面的讨论同样可以清楚地看到，零模 $\eta\dot{x}_c(\tau)$ 对应于反弹的位置，而反弹的位置在式 (2.4.1) 中已经被积分过了。因此，式 (2.4.2) 并不完全正确。实际上我们有
\begin{equation*}
\int \mathrm{d}\tau_1\, K\mathrm{e}^{-S_0} = \frac{\int_{x_c} \mathcal{D}\delta x\, \mathrm{e}^{-S}}{\int_{x=x_0} \mathcal{D}\delta x\, \mathrm{e}^{-S}}\tag{2.4.4}
\end{equation*}
为了计算式 (2.4.4)，我们需要更仔细地定义路径积分。设 $x_n(\tau)$ 是算符 $-m\frac{\mathrm{d}^2}{\mathrm{d}\tau^2} + V''(x_c)$ 的第 $n$ 个归一化本征态，则 $\delta x$ 可以写成 $\delta x = \sum c_n x_n(\tau)$，并且可以定义一个截断的“配分函数”：
\begin{align*}
Z_N(x_c) &\equiv A_N \int \prod_1^N \frac{\mathrm{d}c_n}{\sqrt{2\pi}}\, \mathrm{e}^{-\int \mathrm{d}\tau\, \frac{1}{2}\delta x(-m\frac{\mathrm{d}^2}{\mathrm{d}\tau^2} + V''(x_c(\tau))]\delta x} \\
&= A_N \left(\operatorname{Det}_N\left(-m\frac{\mathrm{d}^2}{\mathrm{d}\tau^2} + V''(x_c)\right)\right)^{-1/2}
\end{align*}
其中 $\operatorname{Det}_N$ 只包括前 $N$ 个本征值的乘积。真正的“配分函数” $Z$ 是 $Z_N$ 在 $N \to \infty$ 下的极限。由于零模 $x_1(\tau)$，上述结果需要修改。我们需要把零模分离出来，得到
\begin{equation*}
Z_N(x_c) = A_N \int \frac{\mathrm{d}c_1}{\sqrt{2\pi}} \left(\operatorname{Det}_N'\left(-m\frac{\mathrm{d}^2}{\mathrm{d}\tau^2} + V''(x_c)\right)\right)^{-1/2}
\end{equation*}
其中 $\operatorname{Det}'$ 表示不包含零本征值。

积分 $\int \mathrm{d}c_1$ 实际上是对反弹位置的积分，也就是 $\int \mathrm{d}\tau_1$。可以求出 $c_1$ 与 $\tau_1$ 通过 $\mathrm{d}\tau_1 = \frac{\mathrm{d}c_1}{\sqrt{S_0/m}}$ 相联系。\footnote{我们注意到
\begin{equation*}
\frac{m}{2}\frac{\mathrm{d}}{\mathrm{d}\tau}(\dot{x}_c)^2 = m\dot{x}_c\ddot{x}_c = \dot{x}_c V'(x_c) = \frac{\mathrm{d}}{\mathrm{d}\tau}V(x_c)
\end{equation*}
因此，若假定 $V(x_0) = 0$，则有 $\frac{m}{2}(\dot{x}_c)^2 = V(x_c)$。由于
\begin{equation*}
S_0 = \int \mathrm{d}\tau\, \Big[\frac{1}{2}m\dot{x}_c^2 + V(x_c)\Big] = \int \mathrm{d}\tau\, m\dot{x}_c^2
\end{equation*}
归一化的零模具有如下形式：
\begin{equation*}
x_1(\tau) = \frac{\dot{x}_c(\tau)}{\sqrt{S_0/m}}
\end{equation*}
于是，$\delta x = c_1 x_1$ 将使反弹产生位移 $\delta\tau = \frac{c_1}{\sqrt{S_0/m}}$。这使我们能够把 $\mathrm{d}c_1$ 换成 $\mathrm{d}\tau_1$。}

因此，
\begin{equation*}
Z_N(x_c) = A_N \int \mathrm{d}\tau_1\, \sqrt{\frac{S_0}{2\pi m}} \left(\operatorname{Det}_N'\left(-m\frac{\mathrm{d}^2}{\mathrm{d}\tau^2} + V''(x_c)\right)\right)^{-1/2}\tag{2.4.5}
\end{equation*}
类似地，对于围绕 $x = x_0$ 的路径积分，我们也可以引入 $Z_N(x_0)$：
\begin{equation*}
Z_N(x_0) = A_N \left(\operatorname{Det}_N\left(-m\frac{\mathrm{d}^2}{\mathrm{d}\tau^2} + V''(x_0)\right)\right)^{-1/2}\tag{2.4.6}
\end{equation*}
现在我们可以写出
\begin{equation*}
\int \mathrm{d}\tau\, K = \lim_{N\to\infty} \frac{Z_N(x_c)}{Z_N(x_0)}
\end{equation*}
由此得到
\begin{align*}
K &= \sqrt{\frac{S_0}{2\pi m}} \left(\frac{\operatorname{Det}\left(-m\frac{\mathrm{d}^2}{\mathrm{d}\tau^2} + V''(x_0)\right)}{\operatorname{Det}'\left(-m\frac{\mathrm{d}^2}{\mathrm{d}\tau^2} + V''(x_c)\right)}\right)^{1/2} \\
&= \sqrt{\frac{S_0}{2\pi}} \left(\frac{\operatorname{Det}\left(-\frac{\mathrm{d}^2}{\mathrm{d}\tau^2} + m^{-1}V''(x_0)\right)}{\operatorname{Det}'\left(-\frac{\mathrm{d}^2}{\mathrm{d}\tau^2} + m^{-1}V''(x_c)\right)}\right)^{1/2}\tag{2.4.7}
\end{align*}
其中用到了这样一个事实：$\operatorname{Det}$ 比 $\operatorname{Det}'$ 多包含一个本征值。同一事实还告诉我们，$K$ 具有 $\tau^{-1}$ 的量纲。

$-\frac{\mathrm{d}^2}{\mathrm{d}\tau^2} + m^{-1}V''(x_0)$ 的所有本征值都是正的。算符 $-\frac{\mathrm{d}^2}{\mathrm{d}\tau^2} + m^{-1}V''(x_c)$ 包含一个零模 $\dot{x}_c(\tau)$，它有一个节点。因此，$-\frac{\mathrm{d}^2}{\mathrm{d}\tau^2} + m^{-1}V''(x_c)$ 有一个、并且只有一个负本征值，正是它使 $K$ 成为虚数。

同一公式 (2.4.7) 也适用于 2.4.1 节中所考虑的隧穿问题。不过，在该情形中零模 $\dot{x}_c$ 没有节点，$-\frac{\mathrm{d}^2}{\mathrm{d}\tau^2} + m^{-1}V''(x_c)$ 没有负本征值。因此，对隧穿问题来说 $K$ 是实数。

我们已经提到，系数 $K$ 来自围绕驻定路径的涨落。从上面的计算我们看到，如果只计入小涨落，$K$ 就可以用算符行列式来表示。Coleman (1985) 讨论过计算行列式之比的若干方法，我在此不再深入。

可能有人已经意识到，眼前这个问题用 WKB 方法很容易解决。路径积分方法实在太繁琐了，人们也许会问，为什么我们还要费心讨论它。进行这一讨论的动机，首先在于它展示了一次典型的路径积分计算中的各个步骤及其微妙之处；其次，也是更重要的一点，同一方法可以用来研究场论中亚稳态的衰变。

\subsection*{问题 2.4.2}
带有衰变亚稳态的最简单场论模型由下式给出：
\begin{equation*}
S = \int \mathrm{d}t\, \mathrm{d}^d\pmb{x}\, \left[\frac{1}{2}\left[(\partial_t m)^2 - v^2(\partial_{\pmb{x}} m)^2\right] - \frac{g}{4}(m^2 - m_0^2)^2 - Bm\right]
\end{equation*}
其中 $m(t, \pmb{x})$ 是一个实场，可以把它看作磁矩的密度。我们将选取单位使速度 $v = 1$。当磁场 $B$ 为零时，系统有两个简并的基态 $m = \pm m_0$，它们是势 $V(m) = \frac{g}{4}(m^2 - m_0^2)^2 + Bm$ 的两个极小点。
\begin{enumerate}
\item 证明：当 $B$ 较小时，系统有一个基态 $m = m_+$ 和一个亚稳态 $m = m_-$。但当 $B$ 超过临界值 $B_c$ 时，亚稳态便不复存在。求 $B_c$ 的值。
\item 写出虚时中的作用量。Coleman 证明，该作用量总有如下形式的反弹解（一条驻定路径）：
\begin{equation*}
m = f\left(\sqrt{(\tau - \tau_0)^2 + (\pmb{x} - \pmb{x}_0)^2}\right) + m_-, \qquad f(\infty) = 0
\end{equation*}
其中 $(\tau_0, \pmb{x}_0)$ 是反弹的位置。
\item 假定在没有反弹时 $\langle m_-|\mathrm{e}^{-TH}|m_-\rangle = \mathrm{e}^{-TE_0}$。把反弹包括进来之后，我们有
\begin{equation*}
\langle m_-|\mathrm{e}^{-TH}|m_-\rangle = \mathrm{e}^{-TE_0} \sum_n \frac{1}{n!} \int \prod_{i=1}^n \mathrm{d}\tau_i\, \mathrm{d}^d\pmb{x}_i\, K^n \mathrm{e}^{-nS_0}
\end{equation*}
其中 $S_0$ 是反弹的作用量，$(\tau_i, \pmb{x}_i)$ 是各个反弹的位置。（注意，现在一个反弹由它在时间与空间两个方面的位置来参数化。）若 $K$ 是虚数，亚稳态 $|m_-\rangle$ 的衰变率是多少？单位体积的衰变率是多少？
\item 设 $B = \frac{B_c}{2}$，估算 $S_0$ 的值。（提示：对 $(\tau, \pmb{x})$ 和 $m$ 重新标度，把 $S$ 改写为 $S = \text{constant} \times S'$，使 $S'$ 中所有系数都是 1 的量级。）（用这一方法也可以估计反弹在时空中的大小。）
\item 设 $B = \frac{B_c}{2}$，估算 $K$ 的值。（提示：考虑当 $S \to aS$ 时 $K$ 如何标度。）
\end{enumerate}

% ------------------------------------------------------------
% 块边界备注：
%   起点：书页 46（ch2_c）末段完整结束（"...The Berry phase vanishes if we
%   choose a common phase for all of the spin states on the path."），
%   书页 47 顶部为空格缩进的新段落（However, for a closed path...），
%   故本块正文直接以该新段开始，无需 %JOIN%，也无 ch2_c 收尾内容可写入
%   %--C-TAIL-- 区（按 assemble.py 约定该区只放"上一文件收尾"）。
%   终点：书页 57 末为止于问题 2.4.2 第 5 问；书页 58 顶部为第 6 问
%   （Discuss for which values of g and m_0 ...），按装配约定归 ch2_e
%   的 %--D-TAIL--，本块未收，ch2_e 请勿重复翻译本块已有的第 1–5 问。
% ------------------------------------------------------------
% 新增术语（ch2_d，待并入术语表"新增术语"区）：
%   instanton 瞬子
%   multi-instanton 多瞬子
%   instanton gas 瞬子气体
%   bounce 反弹
%   tunneling 隧穿
%   tunneling time 隧穿时间
%   barrier 势垒
%   double-well potential 双势阱
%   meta-stable state 亚稳态
%   decay rate 衰变率
%   frustration 阻挫
%   parallel transportation 平行输运
%   geometric phase 几何相位
%   solid angle 立体角
%   direct product 直积
%   spin chain 自旋链
%   spin wave 自旋波
%   ferromagnetic state 铁磁态
%   anti-ferromagnetic state 反铁磁态
%   dispersion relation 色散关系
%   non-perturbative effect 非微扰效应
%   zero mode 零模
%   zero eigenvalue 零本征值
%   node 节点
%   determinant 行列式
%   Hilbert space 希尔伯特空间
%   Newton's equation 牛顿方程
%   magnetic moment 磁矩
%   flux quanta 磁通量子
%   Landau level structure 朗道能级结构
% ------------------------------------------------------------
